%% ========================================================================
%% Sinopsis (Sampul Belakang)
%% ========================================================================

\begin{titlepage}
\thispagestyle{empty}

\begin{tikzpicture}[remember picture, overlay]
    \fill[black!4] ([xshift=1.3cm,yshift=-1.3cm]current page.north west) rectangle
                    ([xshift=-1.3cm,yshift=-4.0cm]current page.north east);
    \draw[line width=0.9pt, black!55]
        ([xshift=1.5cm,yshift=1.5cm]current page.south west) rectangle
        ([xshift=-1.5cm,yshift=-1.5cm]current page.north east);
\end{tikzpicture}

\vspace*{1.6cm}

{\small\sffamily\bfseries S I N O P S I S\par}
\vspace{0.45cm}
{\Large\bfseries Mastering Linux\par}
{\normalsize Pemahaman Sistem Operasi untuk Pemula\par}

\vspace{0.9cm}

\begin{tcolorbox}[
    enhanced,
    colback=white,
    colframe=black!45,
    boxrule=0.5pt,
    arc=1.5pt,
    left=10pt,
    right=10pt,
    top=10pt,
    bottom=10pt
]
Sistem operasi adalah fondasi dari seluruh ekosistem komputasi modern. Tanpa pemahaman yang kuat terhadap cara kerja sistem operasi, seorang praktisi teknologi informasi tidak dapat mengoptimalkan, memvalidasi, atau mengamankan infrastruktur yang dikelolanya.

\vspace{0.5cm}

Buku ini mengajak pembaca memahami Linux dari fondasi hingga praktik administrasi yang digunakan oleh para profesional. Setiap konsep teoritis diperkuat langsung melalui latihan perintah yang terstruktur, sehingga pembaca tidak hanya memahami cara kerja sistem, tetapi juga mampu mengoperasikannya di lingkungan nyata.

\vspace{0.5cm}

Dua belas bab disusun dalam empat bagian: dasar-dasar sistem operasi dan instalasi Linux; manajemen file dan proses; pemrograman Bash shell; serta manajemen sistem lanjutan yang mencakup memori, layanan sistem, paket aplikasi, dan pemulihan sistem. Dengan cakupan yang komprehensif dan gaya penyajian yang terstruktur, buku ini menjadi bekal yang kokoh bagi siapa pun yang ingin membangun kompetensi di bidang administrasi sistem dan DevOps.
\end{tcolorbox}

\vfill

\begin{center}
    {\normalsize\bfseries Polinema Press\par}
    {\small Politeknik Negeri Malang \textbullet{} Malang, Indonesia\par}
\end{center}

\end{titlepage}
